\section*{An exact fixed-axis singularity of the arithmetic trace}
\label{sec:ns-axis-arithmetic-jet}

The angular-momentum transform can see a fixed-axis derivative of the
actual source field directly, without first applying its inverse. The
reason is a specific feature of the construction: every positive-order
azimuthal coefficient has zero axis value. We retain the source amplitude
\(C\), the parameter \(h\), the physical viscosity \(\nu\), all cutoffs and
the original time \(t=1-\tau\). The calculation uses the constructed
field of the preceding actual-source appendix; it does not certify the
entire existence argument of the released manuscript independently.


The source amplitude \(C\) in this appendix is a positive scalar. The
orthogonal operator called \(C\) in the preceding composition remains
exactly \(\mathcal R\Pi\); these two typed uses are stated explicitly
without altering either object. The scalar source parameter \(h\) likewise
retains its original value \(0<h<1/100\); when the solenoidal coefficient
function is needed below it is denoted \(h_{\rm sol}=K_\nu^*u_\nu\).

\subsection*{The axis trace of the fully corrected velocity}

Here and only here let \(q_{\rm src}(z,t)\) denote the scalar similarity
coordinate called \(q\) in the source, not the material label in
the preceding material-flow appendix. Its source coordinates are
\[
 \tau=1-t=q_{\rm src}(1-\eta^2),\qquad
 z=q_{\rm src}^{D}\eta,\qquad
 X=\frac{r^2}{2q_{\rm src}},\qquad
 D=\frac12-h,\quad A=\frac12+h.
\]
In particular at \(z=0\) one has exactly \(\eta=0\) and
\(q_{\rm src}=\tau\).
Source (5.1), pp.~45--46, writes each azimuthal coefficient as
\[
 u_{\theta,n}
   =\frac rC q_{\rm src}^{-1-h+2nh}\varphi_n(X,\eta).
\]
Its Lemma~5.1, p.~47, prescribes
\(\varphi_n(0,\eta)=0\) for every \(n\ge1\).
The radial extensions and moment corrections are supported away from
the unchanged inner interval. Source Proposition~5.5 and its proof,
pp.~60--62, realize this series with cutoffs \(\chi(c_nq_{\rm src})\);
these depend on \(z,t\), not \(r\), and the realized angular component is
\begin{equation}\label{naj:base-series}
 \frac{u_{\theta,B}}r
 =\frac{q_{\rm src}^{-1-h}}C
 \left(\varphi_0(X,\eta)+
   \sum_{n\ge1}\chi(c_nq_{\rm src})
                       q_{\rm src}^{2nh}\varphi_n(X,\eta)\right).
\end{equation}
For every fixed preterminal \(z,t\), this sum is finite in a neighborhood
of that point: \(c_n\to\infty\), while \(q_{\rm src}>0\) and \(\chi\)
has compact support. Thus evaluation at \(X=0\), and the smooth extension
of the quotient \(u_{\theta,B}/r\), require no exchange of a terminal
limit with an infinite sum.

The leading axis datum is not an unknown positive remainder. Source
(B.3), p.~145, and Proposition~B.2, pp.~146--147, give
\[
 \varphi_0(0,\eta)=\varphi_*(\eta),\qquad
 \varphi_*(\eta)=
  \exp\left(\Lambda\int_0^\eta\zeta_*(w)\,dw\right),\qquad
 \varphi_*(0)=1.
\]
Here \(\zeta_*\) is the source's auxiliary axis function, not the
Riemann zeta function. Neither its value nor \(C\) is rescaled.

All annular wave and mean potentials, including the finite
initialization and the later cutoff-summed stages, vanish near the axis
for each \(t<1\): source Proposition~9.9, (9.21), pp.~114--117.
Only finitely many stages are present locally at such a time. Their
Cartesian derivatives therefore vanish on a common neighborhood of the
axis. The base potential is azimuthal and independent of angle, so its
curl has only radial and axial components. Consequently it contributes
no additional azimuthal term to \eqref{naj:base-series}.
Finally the spatial and temporal cutoffs in source (10.4), p.~118,
are identically one near the origin for all sufficiently small
positive \(\tau\). Their derivatives vanish there as well.
These are facts about the actual summed source field, not a
replacement of that field by its leading approximation
\cite[Lemma 5.1; Proposition 5.5; (9.21); (10.4); (B.3)]{openaiNS2026}.

It follows that there is a source-determined \(\tau_1>0\) such that
\begin{equation}\label{naj:axis-star}
 \left.\frac{u_{\theta,*}(r,0,1-\tau)}r\right|_{r=0}
       =\frac1C\tau^{-1-h}\qquad(0<\tau<\tau_1).
\end{equation}
The quotient denotes its proved smooth extension, not division by
zero. For the physical viscosity map
\(u_\nu(x,t)=\sqrt\nu\,u_*(x/\sqrt\nu,t)\), put
\(r_*=r/\sqrt\nu\). Then, for \(r>0\),
\[
 \frac{u_{\theta,\nu}(r,0,t)}r
 =\frac{\sqrt\nu\,u_{\theta,*}(r_*,0,t)}{\sqrt\nu\,r_*}
 =\frac{u_{\theta,*}(r_*,0,t)}{r_*}.
\]
Passing to the smooth axis extension proves the same identity
\eqref{naj:axis-star} for \(u_\nu\). This cancellation follows from the
actual coordinate map; viscosity has not been set to one.

Use the original physical radial coordinate
\(\sigma=r^2/2\) and angular average \( (2\pi)^{-1}\int_0^{2\pi}\).
Let
\[
 B_\nu(\sigma,z,t)=
    \frac1{2\pi}\int_0^{2\pi} r u_{\theta,\nu}(r,\theta,z,t)\,d\theta .
\]
In the axis neighborhood just established the full field is
axisymmetric and \(B_\nu=2\sigma(u_{\theta,\nu}/r)\).
The exact consequence is
\begin{equation}\label{naj:B-axis}
 B_\nu(0,0,1-\tau)=0,\qquad
 \partial_\sigma B_\nu(0,0,1-\tau)=\frac2C\tau^{-1-h}.
\end{equation}
No \(O(\tau^{2h})\) error is present in this first axis derivative:
each positive-order axis value is exactly zero. This does not remove
the actual remainder at the distinct nonzero \(X_{\rm in}\) sample
of the preceding actual-source appendix.

\subsection*{Complete endpoint summation formula and remainder}

We prove the needed endpoint summation formula explicitly. Let
\(B_m(v)\) be the Bernoulli polynomials defined by \(B_0=1\),
\(B_m'=mB_{m-1}\), and \(\int_0^1B_m(v)\,dv=0\) for \(m\ge1\);
write \(B_m=B_m(0)\). In particular \(B_1(v)=v-1/2\).
For \(m\ge2\) integration of \(B_m'\) shows \(B_m(1)=B_m(0)\).
Uniqueness of this defining recursion proves
\(B_m(1-v)=(-1)^mB_m(v)\), so \(B_m(0)=0\) for odd \(m\ge3\).
Let \(\widetilde B_m(v)=B_m(\{v\})\); the values at the integers
are immaterial inside integrals.

Suppose \(f\in C^{2M}([0,\infty))\) and its derivatives through order \(2M\) are integrable,
have the indicated one-sided values at zero, and tend to zero at infinity
through order \(2M-1\). Summing integration by parts over the intervals
\([nx,(n+1)x]\) gives
\[
 \sum_{n\ge1}f(nx)
 =x^{-1}\int_0^\infty f(v)\,dv-\frac{f(0)}2
       +\int_0^\infty\widetilde B_1(v/x)f'(v)\,dv .
\]
For example, the integral on one interval equals the half-sum of its
endpoint values minus \(x^{-1}\) times its integral. For every subsequent
integration the identity
\(\frac{d}{dv}\widetilde B_m(v/x)=
m x^{-1}\widetilde B_{m-1}(v/x)\) holds between endpoints.
The endpoint values match for \(m\ge2\), and the odd endpoint constants
vanish as just proved. Repeating these integrations gives
\begin{align}
 \sum_{n\ge1}f(nx)
 &=x^{-1}\int_0^\infty f(v)\,dv-\frac{f(0)}2
   -\sum_{k=1}^M\frac{B_{2k}}{(2k)!}
           f^{(2k-1)}(0)x^{2k-1}+R_M(f,x),\label{eq:nh-EM}\\
 |R_M(f,x)|&\le
 \frac{\|B_{2M}\|_{L^\infty[0,1]}}{(2M)!}
 x^{2M-1}\int_0^\infty|f^{(2M)}(v)|\,dv .\notag
\end{align}
The remainder is the negative of
\(x^{2M-1}(2M)!^{-1}\int\widetilde B_{2M}(v/x)f^{(2M)}(v)\,dv\);
its displayed bound justifies the infinite endpoint limit as well.


The recursion gives \(B_2(v)=v^2-v+1/6\), so \(B_2=1/6\).
For the compact smooth inputs below all integrability and boundary
hypotheses of this formula hold at every finite order.

\subsection*{The arithmetic operator on every endpoint derivative}

We prove the endpoint map on its full smooth compactly supported
domain. Let \(a\in C^\infty([0,\infty))\) have compact support and
\(a(0)=0\). For the same fixed \(h\) define
\[
 \phi_a(x)=a(x)-2^h a(2x),\qquad
 \psi_a(x)=\phi_a(x)-\tfrac12\phi_a(x/2),\qquad
 (\mathcal A_h a)(x)=\sum_{n\ge1}\psi_a(nx)\quad(x>0).
\]
This is exactly the arithmetic operator used on the nonlinear fluxes.
There is no radius change hidden in its argument \(x\).
Every sum and its derivatives are locally finite for \(x>0\).
Substitution in the two integrals proves
\(\int_0^\infty\psi_a=0\), and \(\psi_a(0)=0\).
The full endpoint derivatives of the input are
\begin{equation}\label{naj:psi-jet}
 \psi_a^{(j)}(0)
  =(1-2^{-j-1})(1-2^{h+j})a^{(j)}(0)\qquad(j\ge0).
\end{equation}

For clarity the passage from the sum to actual endpoint derivatives
is justified here, rather than differentiating a formal asymptotic
series. We use the Euler--Maclaurin identity proved with its remainder
in \eqref{eq:nh-EM}; its classical provenance is
\href{https://dlmf.nist.gov/2.10.E1}{DLMF, \S2.10(i), (2.10.1)}. With
\(f_j(v)=v^j\psi_a^{(j)}(v)\), direct differentiation on \(x>0\)
gives
\[
 x^j(\mathcal A_ha)^{(j)}(x)=\sum_{n\ge1}f_j(nx).
\]
All these \(f_j\) are smooth and compactly supported. Integration by
parts \(j\) times gives
\(\int f_j=(-1)^j j!\int\psi_a=0\); the factors \(v^j\) make
each zero-end boundary term vanish. Their axis derivatives are
\[
 f_j^{(r)}(0)=0\quad(r<j),\qquad
 f_j^{(r)}(0)=\frac{r!}{(r-j)!}\psi_a^{(r)}(0)\quad(r\ge j).
\]
Choose the Euler--Maclaurin order \(M\) so that \(2M-1>j\).
After division by \(x^j\), its remainder is bounded by
\[
 \frac{\|B_{2M}\|_\infty}{(2M)!}\,
 x^{2M-1-j}\|f_j^{(2M)}\|_1,
\]
and hence tends to zero. Terms below degree \(j\) vanish, so every
differentiated series has a finite limit at zero. The only constant
term for \(j\ge1\) is the Bernoulli term with \(2k-1=j\); for even
\(j\ge2\) it is absent. The fundamental theorem of calculus on
\([\epsilon,x]\), followed by \(\epsilon\downarrow0\), now identifies
these derivative limits with successive derivatives of the extended
function. Boundedness of the next derivative justifies the passage
under each integral. Thus \(\mathcal A_ha\) is smooth at zero, and
\begin{equation}\label{naj:all-jets}
 \begin{split}
 (\mathcal A_ha)(0)&=0,\\
 (\mathcal A_ha)^{(j)}(0)&=
 -\frac{B_{j+1}}{j+1}
 (1-2^{-j-1})(1-2^{h+j})a^{(j)}(0),\qquad j\ge1.
 \end{split}
\end{equation}
The Bernoulli convention is \(B_2=1/6\) and odd \(B_{j+1}=0\)
for even \(j\ge2\).
In particular
\begin{equation}\label{naj:first-map}
 (\mathcal A_ha)'(0)=
       \frac{2^{h+1}-1}{16}\,a'(0).
\end{equation}
The coefficient is strictly positive for the source \(h>0\).

The preceding proofs apply also after any fixed \(z,t\) derivative of
a smooth compactly supported parameter family with common support on
compact preterminal parameter sets. Every bound then holds uniformly
on those sets, with constants depending on the derivative order.
This proves compatibility of the endpoint map with those parameter
derivatives. It supplies no uniform estimate as \(t\uparrow1\).

\subsection*{A directly visible singularity at a fixed arithmetic coordinate}

Apply this same operator to the actual compactly supported
\(a(\sigma)=B_\nu(\sigma,z,t)\). Write its arithmetic image as
\[
 \mathcal B_\nu(x,z,t)=\mathcal A_hB_\nu(x,z,t).
\]
The domain and all nonlinear fluxes are those already given in the
preceding angular-momentum section. Equations \eqref{naj:B-axis} and
\eqref{naj:first-map} now prove the exact formula
\begin{equation}\label{naj:exact-blowup}
 \boxed{\displaystyle
 \partial_x\mathcal B_\nu(0,0,1-\tau)
    =\frac{2^{h+1}-1}{8C}\,\tau^{-1-h},
       \qquad 0<\tau<\tau_1,\quad \nu>0 .}
\end{equation}
Here \(x=0,z=0\) are fixed coordinates, not a shrinking sampling path,
and \(\mathcal B_\nu(0,z,t)=0\) for every preterminal \(t\).
Thus the value is zero while its first \(x\)-derivative diverges
positively. For any fixed \(\epsilon>0\), the \(C^1([0,\epsilon])\)
seminorm is at least the displayed endpoint derivative. In particular
no extension continuous at \(t=1\) with values in
\(C^1([0,\epsilon])\) can agree with this family.
This assertion follows from an exact derivative of the original
arithmetic sum, not merely a lower bound obtained after its inverse.

For every integer \(m\ge0\), differentiation of the exact equality gives
\[
 \partial_t^m\partial_x\mathcal B_\nu(0,0,t)
  =\frac{2^{h+1}-1}{8C}(1+h)_m(1-t)^{-1-h-m},
 \qquad 1-\tau_1<t<1 .
\]
Indeed \(d_t(1-t)^{-\alpha}=\alpha(1-t)^{-\alpha-1}\).
These are equalities on an open preterminal interval. They make no
assignment of finite terminal derivatives.

The projection onto \(B_\nu\) has not discarded the other velocity
data in the full correspondence. Its exact smooth section and kernel
are
\[
 R B_\nu=\frac{B_\nu(r^2/2,z,t)}{r^2}(-x_2,x_1,0),\qquad
 w_\nu=u_\nu-RB_\nu,\qquad \langle r(w_\nu)_\theta\rangle=0,
\]
where \(B_\nu/\sigma\) has its smooth integral extension at zero.
The preceding section proves the inverse
\((B_\nu,w_\nu)\mapsto RB_\nu+w_\nu\) and retains every covariance
in both momentum fluxes. The signal \eqref{naj:exact-blowup} concerns
the \(B_\nu\) component of that complete map. It does not replace the
full Navier--Stokes equation by a closed scalar equation for \(B_\nu\).

\subsection*{The exact Mellin pole carrying this axis signal}

For fixed preterminal parameters, put
\[
 \widehat B_\nu(s)=\int_0^\infty
                  B_\nu(\sigma,z,t)\sigma^{s-1}\,d\sigma,\qquad
 \mathfrak B_\nu(s)=\int_0^\infty
                  \mathcal B_\nu(x,z,t)x^{s-1}\,dx .
\]
Both integrals are initially holomorphic on \(\Re s>-1\).
The full transform calculation in the preceding section gives
\[
 \mathfrak B_\nu(s)=W_h(s)\widehat B_\nu(s),\qquad
 W_h(s)=\zeta(s)(1-2^{s-1})(1-2^{h-s}).
\]
The pole of \(\zeta\) at \(1\) is canceled by the displayed first
dilation factor, making \(W_h\) entire.

Here is the continuation needed at \(s=-1\). For any fixed \(b>0\),
split the integral at \(b\) and subtract the linear Taylor term on
\([0,b]\). Since \(B_\nu(0)=0\), the exact formula is
\[
 \widehat B_\nu(s)=
 \int_0^b [B_\nu(\sigma)-B_{\nu,\sigma}(0)\sigma]\sigma^{s-1}d\sigma
 +\int_b^\infty B_\nu(\sigma)\sigma^{s-1}d\sigma
 +\frac{B_{\nu,\sigma}(0)b^{s+1}}{s+1}.
\]
Taylor's integral remainder is \(O(\sigma^2)\); the first integral
is holomorphic for \(\Re s>-2\), including after any complex
\(s\) derivative on a compact subset, because the additional powers
of \(\log\sigma\) remain integrable. The second is entire by compact
support away from zero. Hence
\(\operatorname{Res}_{s=-1}\widehat B_\nu=B_{\nu,\sigma}(0)\).
The identical argument applies to the proved smooth
\(\mathcal B_\nu\), so its residue is \(\partial_x\mathcal B_\nu(0)\).


We also derive the required zeta value from the same finite-endpoint
calculation, so that this value does not stand in for a proof.
For \(\Re s>1\), apply the cell integrations on \([1,N]\) to
\(f(v)=v^{-s}\), include the first endpoint, and let \(N\to\infty\).
Since \(f^{(k)}(v)=(-1)^k(s)_k v^{-s-k}\), the exact result is
\begin{align*}
 \zeta(s)={}&\frac1{s-1}+\frac12+
 \sum_{k=1}^{M}\frac{B_{2k}}{(2k)!}(s)_{2k-1}\\
 &-\frac{(s)_{2M}}{(2M)!}
       \int_1^\infty \widetilde B_{2M}(v)v^{-s-2M}\,dv,
 \qquad (s)_k=\prod_{j=0}^{k-1}(s+j),\quad (s)_0=1.
\end{align*}
The integral is holomorphic on \(\Re s>1-2M\): on every compact subset,
its integrand and every complex derivative are bounded by a constant
times \(v^{-1-\delta}(1+|\log v|^j)\) for some \(\delta>0\), an
integrable function. This follows from boundedness of the periodic
Bernoulli polynomial. Identity on the initial half-plane identifies the
meromorphic continuation. Increasing \(M\) covers the whole plane;
the formulas agree on their overlaps by the identity theorem, and the
only pole is the displayed pole at \(s=1\) with residue one.
At \(s=-1\), use \(M=2\). The products \((s)_3\) and \((s)_4\)
vanish there, while the remainder integral is holomorphic there. Thus
\[
 \zeta(-1)=\frac1{-1-1}+\frac12+
                \frac{B_2}{2!}(-1)
          =-\frac12+\frac12-\frac1{12}=-\frac1{12}.
\]
This also proves directly the claimed cancellation of the sole zeta pole
by \(1-2^{s-1}\), since its derivative at one is \(-\log2\).

The classical value \(\zeta(-1)=-1/12\)
\href{https://dlmf.nist.gov/25.6}{DLMF, \S25.6(i)} gives
\[
 W_h(-1)=-\frac1{12}\left(1-\frac14\right)(1-2^{h+1})
         =\frac{2^{h+1}-1}{16}>0.
\]
Equality on \(\Re s>-1\) therefore continues to the meromorphic
functions just constructed and agrees with the direct endpoint proof.
At \(z=0,t=1-\tau\) the pole is genuine and its residue is exactly
\begin{equation}\label{naj:residue}
 \operatorname{Res}_{s=-1}\mathfrak B_\nu(s,0,1-\tau)
     =\frac{2^{h+1}-1}{8C}\tau^{-1-h}.
\end{equation}
There is no zero-factor compensation at this pole, since \(W_h(-1)>0\).
The sign of \(\zeta(-1)\) alone is not the sign of the arithmetic
residue: both dilation factors have been evaluated explicitly.

This identifies what the construction propagates in this channel:
an exact axis derivative, equivalently a growing residue at the fixed
Mellin point \(-1\). The identity retains the full
\(\widehat B_\nu\), its flux equations, and its inverse maps.
The nontrivial zeros of \(\zeta\) remain the zero factors specified
in the full product formula; no off-critical zero is inferred from
this axis pole or from its positive residue.
